\subsection*{一步一步建立直观}

\paragraph{1. 先写出递推，而非直接画无穷高}

把有限幂塔看作迭代 $x_{k+1}=b^{x_k}$。右图的蜘蛛网先竖直到达函数曲线，再横着回到对角线，把输出作为下一次输入。若序列真的收敛且指数函数连续，极限必满足 $x=b^x$；这一步给出必要条件，还没有保证收敛。



\paragraph{2. 交替接近也可能收敛}

当底数小于一时，指数函数递减，迭代可能在不动点两侧往返。检查附近导数的绝对值是否小于一，有助于理解小误差是被缩小还是放大。图中几次迭代看似稳定只是线索；局部收缩还需结合进入该邻域的条件。



\paragraph{3. 不动点方程不能替代收敛证明}

将底数提高后，轨迹越走越大，最终可能超出显示范围。程序会明确报告溢出范围，而不会把屏幕边缘当成极限。无穷表达式若没有先定义为一个收敛过程，就不能随意把“顶端的无穷幂塔”替换为自身并当作求值完成。



\subsection*{定义、命题与证明}

\paragraph{定义：幂塔的有限迭代}

给定底数 $b>0$ 和起点 $x_0$，定义 $x_{k+1}=b^{x_k}$。所谓无限幂塔的取值，需要先说明这个序列的极限是否存在。

\paragraph{命题：收敛的迭代必到达不动点}

若 $x_k$ 收敛到有限实数 $L$，则 $L=b^L$。该方程只是收敛的必要条件。

\begin{proof}

收敛序列平移一项仍趋于 $L$，所以 $x_{k+1}\to L$。函数 $x\mapsto b^x$ 在实数上连续，故 $b^{x_k}\to b^L$。

迭代两边逐项相等，极限的唯一性给出 $L=b^L$。证明从“序列收敛”开始，不能倒过来用同一个推理宣称只要方程有解，所有初值的迭代就会收敛。

\end{proof}
